\section*{Capitulo 10. Repaso acumulativo, simulacros y reto final}
\addcontentsline{toc}{section}{Capitulo 10. Repaso acumulativo, simulacros y reto final}

\begin{shortanswer}{R10.S01 diagnostico}
\begin{enumerate}[label=\textbf{PX-R10.S01-\arabic*:}, leftmargin=*]
  \item \((1,\infty)\)
  \item \(x=4\)
  \item \(x=\dfrac{\pi}{3}\) o \(x=\dfrac{5\pi}{3}\)
  \item \(y=-x-1\)
\end{enumerate}
\end{shortanswer}

\begin{shortanswer}{R10.S02 hoja mixta}
\begin{enumerate}[label=\textbf{PX-R10.S02-\arabic*:}, leftmargin=*]
  \item \(y=-\dfrac{1}{2}x+\dfrac{5}{2}\)
  \item \(4\)
  \item El cuadrado de lado \(5\), con area \(25\)
  \item \((2,3)\) y \((3,2)\)
\end{enumerate}
\end{shortanswer}

\begin{shortanswer}{R10.S03 eleccion de metodo}
\begin{enumerate}[label=\textbf{PX-R10.S03-\arabic*:}, leftmargin=*]
  \item \(x=3\)
  \item Asintota vertical \(x=1\); asintota oblicua \(y=x+1\)
  \item Crece en \((-\infty,-1)\cup(1,\infty)\), decrece en \((-1,1)\), maximo local
  \((-1,2)\), minimo local \((1,-2)\)
\end{enumerate}
\end{shortanswer}

\begin{shortanswer}{R10.S04 analisis de errores}
\begin{enumerate}[label=\textbf{PX-R10.S04-\arabic*:}, leftmargin=*]
  \item Falta el caso \(2x-1=-5\); la solucion correcta es \(x=3\) o \(x=-2\)
  \item Hay que multiplicar por el conjugado; el resultado correcto es
  \(\dfrac{\sqrt{3}+1}{2}\)
  \item La pendiente normal es la opuesta inversa, \(-\dfrac{1}{2}\); la normal es
  \(y=-\dfrac{1}{2}x+\dfrac{3}{2}\)
\end{enumerate}
\end{shortanswer}

\begin{shortanswer}{R10.S05 simulacro I}
\begin{enumerate}[label=\textbf{SM-R10.S05-\arabic*:}, leftmargin=*]
  \item \((-\infty,-1]\cup(2,\infty)\)
  \item \(x=\dfrac{\pi}{3}\) o \(x=\dfrac{2\pi}{3}\)
  \item \(a=2\)
  \item \(y=4x-3\)
\end{enumerate}
\end{shortanswer}

\begin{shortanswer}{R10.S06 simulacro II}
\begin{enumerate}[label=\textbf{SM-R10.S06-\arabic*:}, leftmargin=*]
  \item \(x=1+\sqrt{2}\) o \(x=1-\sqrt{2}\)
  \item \(\dfrac{6}{\sqrt{5}}=\dfrac{6\sqrt{5}}{5}\)
  \item \(y=x-1\)
  \item Minimo \(4\), alcanzado en \(x=2\)
\end{enumerate}
\end{shortanswer}

\begin{shortanswer}{[CH-R10-01]}
Dominio \(\R\setminus\{0\}\), asintotas \(x=0\) e \(y=x\), crecimiento en
\((-\infty,-2)\cup(2,\infty)\), decrecimiento en \((-2,0)\cup(0,2)\), maximo local
\((-2,-4)\), minimo local \((2,4)\), tangente en \(x=2\): \(y=4\).
\end{shortanswer}
